% =====================================================================
\section{Asymmetric defence}\label{sec:asym}
\figp[0.88\textwidth]{fig_v5.png}{\label{fig:v5}Four results of version~5. \textbf{a} --- asymmetric defence: the measured level lies on the derivation $w(r)$, the ordering agreed exactly. The range is anchored by the two extreme points (circled); the three middle ones are predicted. \textbf{b} --- burstiness of the load: with the mean and the variance pinned, the collapse rate falls by 40\,\%, whereas the fraction of cells that were damaged saturates and sees nothing; the choice of the measured quantity decided the outcome of the rule. \textbf{c} --- the width of the rescued band against the damaged fraction: two scans give opposite signs, so the band is not a function of $D$. \textbf{d} --- the cost per row doubles for the one-operation-at-a-time computation once the working set no longer fits in cache, and does not change at all for the fused kernel.}

The second stage of the defence had all along been written with a single constant: $\dot s=(m-s)/\tau_s$, 63 hours both on the rise and on the decay. Biologically these are two different processes --- induction of synthesis and degradation of protein --- and their equality was an assumption we had simply never tested.\tD\ Version~5 separates them: $\tau_{\uparrow}$ on the rise, $\tau_{\downarrow}$ on the decay, with the ratio $r=\tau_{\downarrow}/\tau_{\uparrow}$.

\textbf{A derivation yielding a prediction before any computation.} Under a daily cycle the input $m$ is two-level to good accuracy: the high level holds for a fraction $f$ of the day, the low level for the rest. For a leaky integrator with different constants, when both are much larger than a day, the steady mean stands between the levels at the point
\[w(r)=\frac{f\,r}{f\,r+1-f},\]
that is, $s_{\rm mean}=m_{\rm lo}+w(r)\,(m_{\rm hi}-m_{\rm lo})$. The asymmetry turns the integrator into a \emph{peak detector} for $r>1$ and into a \emph{trough detector} for $r<1$. The fraction $f$ is taken from the run (it equals the measured fraction of waking, 0.705) rather than assigned, so the prediction contains not a single fitted number.

\textbf{What came out.} At $r=0.33/0.5/1/2/3$ the derivation gives $w=0.441/0.544/0.705/0.827/0.878$, while the measured level of the defence is $0.0575/0.0644/0.0792/0.0958/0.1065$. The ordering agreed \emph{exactly}: rank association $+1$ at an exact permutation $p=0.0083$ (fig.~\ref{fig:v5}a). The consequence for collapses was confirmed as well: the fraction of cells that were damaged falls from 0.486 to 0.457, the trend test gives $z=1.67$, $p=0.048$, and the episode duration shortens from 7.74 to 6.88 days.

\textbf{The conclusion that changes the status of the earlier model.} The equality of the rise and decay constants was not a neutral simplification but a \emph{choice of regime}. Different $r$ give a different mean defence under one and the same instantaneous equations, and the direction of the effect is predictable from two lines of algebra. So either the equality must be justified by measurement, or $r$ must remain a parameter of the model on the same footing as $\tau_s$.

\begin{note}
The quantitative part of the rule remained unverifiable, and the reason lies in the \emph{construction of the rule itself}, not in a shortage of sets. The equivalence tolerance $3\,\mathrm{SE}=0.0241$ and half the range of the curve, $0.0245$, almost coincided, so the gap between the threshold and the alternative is 0.0004 at an MDE of 0.02. No number of sets cures this: 29 times more would be needed. What cures it is leverage --- widening the scan of $r$ from 0.1 to 10, which increases the range of the curve 1.7-fold. This is the price recorded honestly: we know \emph{what} we did not verify and \emph{how much} it costs.
\end{note}
